\chapter{Fully Lazy Lambda-lifting}
%\vspace{2cm}
\vspace{1cm}

As we discussed in Chapter~11 our implementations support lazy evaluation.
However, there is one major way in which an implementation based on
lambda-lifting can be made still lazier than the version we have described so
far. The purpose of this chapter is to describe the opportunity and the
modifications required to exploit it.

\section{Full Laziness}

As we remarked in Section~12.1.3, a straightforward implementation of the
template-instantiation procedure risks constructing \textit{multiple instances of the
    same expression}, rather than sharing a single copy of them. This wastes space
because each instance occupies separate storage, and it wastes time because
the instances will be reduced separately. This waste can be arbitrarily large;
for example, the duplicated instances might each separately perform some
large calculation.

The loss of sharing can best be seen using an example. Consider the
function
\begin{mlcoded}
    f = \tlb{y}+ y (sqrt 4)
\end{mlcoded}
Whenever this function is applied to an argument we will slavishly construct a
new instance of the subexpression \ml{(sqrt 4)} in its body, despite the fact that all
instances of the \ml{(sqrt 4)} reduce to \ml{2}. It would be better not to construct a new
instance of such constant subexpressions, but to share a single instance
instead. This can do no harm, since the constant subexpression does not
contain any occurrences of the formal parameter, and hence its value cannot
change between one application and another.

It looks as if these constant subexpressions could be spotted and marked by
a compiler, but they can be generated `on the fly'. Consider the Miranda
program

\combinatorbox{
    \ml{f = g 4} \\
    \indent\ml{g x y = y + (sqrt x)}
}{
    \ml{(f 1) + (f 2)}
}

\noindent
This compiles to the lambda expression:

\combinatorbox{
    \begin{mlcoded}
    letrec f = g 4 \\
    \phantom{letrec} g = \tlb{x}\tlb{y}+ y (sqrt x) \\
    in + (f 1) (f 2)
    \end{mlcoded}
    \vspace{-0.5\baselineskip}
}{}

\noindent
Therefore, when evaluating the expression, we get

\vs
\begin{tikzpicture}[
    bullet/.style={circle, fill=black, inner sep=1.5pt},
    arr/.style={-{Stealth[length=2.5mm,width=1.5mm]}, thick},
    node distance=0pt,
    every node/.style={font=\sffamily\footnotesize}
    ]

    \def\rs{0.9}
    \def\as{0.55}

    \pgfmathsetmacro\rowa{0}
    \pgfmathsetmacro\rowb{-1*\rs}
    \pgfmathsetmacro\rowc{-1*\rs-\as}
    \pgfmathsetmacro\rowd{-2*\rs-\as}
    \pgfmathsetmacro\rowe{-2*\rs-2*\as}
    \pgfmathsetmacro\rowf{-3*\rs-2*\as}
    \pgfmathsetmacro\rowg{-3*\rs-3*\as}
    \pgfmathsetmacro\rowh{-4*\rs-3*\as}
    \pgfmathsetmacro\rowi{-4*\rs-4*\as}
    \pgfmathsetmacro\rowj{-5*\rs-4*\as}

    \def\tgtx{4.8}

    % Row 1: + (f 1) (f 2)  [no arrow]
    \node[anchor=north west, inner sep=0pt] at (-1.2,\rowa)
    {\phantom{$\rightarrow$ } \ml{+ (f\quad \ \ 1) (f\ \ \ 2)}};

    % Row 2: -> + (• 1) (• 2)   --> ((λx.λy.+ y (sqrt x)) 4)
    \node[anchor=west] at (-1.2,\rowb) {$\rightarrow$ \ml{+ (}};
    \node[bullet] (b2a) at (0.25,\rowb) {};
    \node[anchor=west] at (0.35,\rowb) {\ml{1) (}};
    \node[bullet] (b2b) at (1.15,\rowb) {};
    \node[anchor=west] at (1.25,\rowb) {\ml{2)}};
    \node[anchor=west] (r2tgt) at (\tgtx,\rowc) {\ml{((\tlb{x}\tlb{y}+ y (sqrt x)) 4)}};
    \draw (b2a.south) -- (0.25,\rowc);
    \draw[arr] (0.25,\rowc) -- (r2tgt.west);
    \draw (b2b.south) -- (1.15,\rowc);
    \draw (1.15,\rowc) -- (\tgtx,\rowc);

    % Row 3: -> + (• 1) (• 2)   --> (λy.+ y (sqrt 4))
    \node[anchor=west] at (-1.2,\rowd) {$\rightarrow$ \ml{+ (}};
    \node[bullet] (b3a) at (0.25,\rowd) {};
    \node[anchor=west] at (0.35,\rowd) {\ml{1) (}};
    \node[bullet] (b3b) at (1.15,\rowd) {};
    \node[anchor=west] at (1.25,\rowd) {\ml{2)}};
    \node[anchor=west] (r3tgt) at (\tgtx,\rowe) {\ml{(\tlb{y}+ y (sqrt 4))}};
    \draw (b3a.south) -- (0.25,\rowe);
    \draw[arr] (0.25,\rowe) -- (r3tgt.west);
    \draw (b3b.south) -- (1.15,\rowe);
    \draw (1.15,\rowe) -- (\tgtx,\rowe);

    % Row 4: -> + (• 1) (+ 2 (sqrt 4))   --> (λy.+ y (sqrt 4))
    \node[anchor=west] at (-1.2,\rowf) {$\rightarrow$ \ml{+ (}};
    \node[bullet] (b4a) at (0.25,\rowf) {};
    \node[anchor=west] at (0.35,\rowf) {\ml{1) (+ 2 (sqrt 4))}};
    \node[anchor=west] (r4tgt) at (\tgtx,\rowg) {\ml{(\tlb{y}+ y (sqrt 4))}};
    \draw (b4a.south) -- (0.25,\rowg);
    \draw[arr] (0.25,\rowg) -- (r4tgt.west);

    % Row 5: -> + (• 1) 4   --> (λy.+ y (sqrt 4))
    \node[anchor=west] at (-1.2,\rowh) {$\rightarrow$ \ml{+ (}};
    \node[bullet] (b5a) at (0.25,\rowh) {};
    \node[anchor=west] at (0.35,\rowh) {\ml{1) 4}};
    \node[anchor=west] (r5tgt) at (\tgtx,\rowi) {\ml{(\tlb{y}+ y (sqrt 4))}};
    \draw (b5a.south) -- (0.25,\rowi);
    \draw[arr] (0.25,\rowi) -- (r5tgt.west);

    % Rows 6–8: plain lines, lumped into one node
    \node[anchor=north west, align=left, inner sep=0pt] at (-1.2,\rowj) {%
        $\rightarrow$ \ml{+ (+ 1 (sqrt 4)) 4}\\[2pt]
        $\rightarrow$ \ml{+ 3 4}\\[2pt]
        $\rightarrow$ \ml{7}%
    };

\end{tikzpicture}
\vs

\noindent
The crucial point is that the \ml{(sqrt 4)} is evaluated twice, because a fresh
instance of it is made each time the \ml{\tl y} is applied. The reason for this is that it is
a \textit{dynamically created} constant subexpression of the \ml{\tl y} abstraction.

Not surprisingly, just the same problem occurs with supercombinators. Our
example compiles to

\combinatorbox{
    \ml{\$g x y = + y (sqrt x)} \\
    \indent\ml{\$f = \$g 4} \\
    \indent\ml{\$Prog = + (\$f 1) (\$f 2)}
}{
    \ml{\$Prog}
}

\noindent
The reduction proceeds as follows:

\vs
\begin{tikzpicture}[
    bullet/.style={circle, fill=black, inner sep=1.5pt},
    arr/.style={-{Stealth[length=2.5mm,width=1.5mm]}, thick},
    node distance=0pt,
    every node/.style={font=\sffamily\footnotesize}
    ]

    \def\rs{0.9}
    \def\as{0.55}

    \pgfmathsetmacro\rowa{0}
    \pgfmathsetmacro\rowb{-1*\rs}
    \pgfmathsetmacro\rowc{-1*\rs-\as}
    \pgfmathsetmacro\rowd{-2*\rs-\as}
    \pgfmathsetmacro\rowe{-2*\rs-2*\as}
    \pgfmathsetmacro\rowf{-3*\rs-2*\as}
    \pgfmathsetmacro\rowg{-3*\rs-3*\as}
    \pgfmathsetmacro\rowh{-4*\rs-3*\as}

    \def\tgtx{4.0}

    % Row 1: $Prog  [no arrow]
    \node[anchor=north west, inner sep=0pt] at (-1.2,\rowa)
    {\phantom{$\rightarrow$ } \ml{\$Prog}};

    % Row 2: -> + (• 1) (• 2)   --> ($g 4)
    \node[anchor=west] at (-1.2,\rowb) {$\rightarrow$ \ml{+ (}};
    \node[bullet] (b2a) at (0.25,\rowb) {};
    \node[anchor=west] at (0.35,\rowb) {\ml{1) (}};
    \node[bullet] (b2b) at (1.15,\rowb) {};
    \node[anchor=west] at (1.25,\rowb) {\ml{2)}};
    \node[anchor=west] (r2tgt) at (\tgtx,\rowc) {\ml{(\$g 4)}};
    \draw (b2a.south) -- (0.25,\rowc);
    \draw[arr] (0.25,\rowc) -- (r2tgt.west);
    \draw (b2b.south) -- (1.15,\rowc);
    \draw (1.15,\rowc) -- (\tgtx,\rowc);

    % Row 3: -> + (• 1) (+ 2 (sqrt 4))   --> ($g 4)
    \node[anchor=west] at (-1.2,\rowd) {$\rightarrow$ \ml{+ (}};
    \node[bullet] (b3a) at (0.25,\rowd) {};
    \node[anchor=west] at (0.35,\rowd) {\ml{1) (+ 2 (sqrt 4))}};
    \node[anchor=west] (r3tgt) at (\tgtx,\rowe) {\ml{(\$g 4)}};
    \draw (b3a.south) -- (0.25,\rowe);
    \draw[arr] (0.25,\rowe) -- (r3tgt.west);

    % Row 4: -> + (• 1) 4   --> ($g 4)
    \node[anchor=west] at (-1.2,\rowf) {$\rightarrow$ \ml{+ (}};
    \node[bullet] (b4a) at (0.25,\rowf) {};
    \node[anchor=west] at (0.35,\rowf) {\ml{1) 4}};
    \node[anchor=west] (r4tgt) at (\tgtx,\rowg) {\ml{(\$g 4)}};
    \draw (b4a.south) -- (0.25,\rowg);
    \draw[arr] (0.25,\rowg) -- (r4tgt.west);

    % Rows 5–7: plain lines
    \node[anchor=north west, align=left, inner sep=0pt] at (-1.2,\rowh) {%
        $\rightarrow$ \ml{+ (+ 1 (sqrt 4)) 4}\\[2pt]
        $\rightarrow$ \ml{+ 3 4}\\[2pt]
        $\rightarrow$ \ml{7}%
    };

\end{tikzpicture}
\vs

\noindent
Again we see that the \ml{(sqrt 4)} has been evaluated twice.

To be as lazy as possible we would like to share even these dynamically
created constant expressions. Specifically, the effect we want to achieve is that
every expression is evaluated \textit{at most once} after the variables in it have been
bound. This is called \textit{full laziness}. It corresponds closely to an optimization
sometimes performed by conventional compilers on loops, in which
expressions not involving the loop variable (i.e.\ free expressions) are moved
out of the loop so that they are not repeatedly evaluated.

\section{Maximal Free Expressions}

The problem we have discovered is that laziness can be lost if we instantiate
too much of the body of a lambda abstraction.

Which parts should not be instantiated? The parts of the body that should
not be instantiated are those subexpressions which contain no (free)
occurrences of the formal parameter, because if the formal parameter does
not occur then the value of the subexpression will be the same between all
instances, and hence may be shared. To formalize this we need a new
definition.

\definitionbox{
    \vspace{-0.8\baselineskip}
}{
    A subexpression \ml{E} of a lambda abstraction \ml{L} is \textit{free} in \ml{L} if all variables in \ml{E}
    are free in \ml{L}. A \textit{maximal free expression} or \textit{MFE} of \ml{L} is a free expression
    which is not a proper subexpression of another free expression of \ml{L}. (\ml{E} is a
    proper subexpression of \ml{F} if and only if \ml{E} is a subexpression of \ml{F} and
    \ml{E} $\neq$ \ml{F}.)
}

\noindent
\textit{Examples}\\
In the following lambda abstractions the maximal free expressions of the \tl x
abstractions are underlined.
\vs

\begin{compactenum}[\quad(1)]
    \item \ml{(\tlb{x}\underline{sqrt} x)}
    \item \ml{(\tlb{x}x \underline{(sqrt 4)})}
    \item \ml{(\tlb{y}\tlb{x}$\pm$ x \underline{($\ast$ y y)})}
    \item \ml{(\tlb{y}\tlb{x}+ \underline{($\ast$ y y)} x)}
    \item \ml{(\tlb{x}\underline{(\tlb{x}x)} x)} (here the \ml{(\tlb{x}x)} is free despite the name clash)
\end{compactenum}
\vs

\noindent
To achieve full laziness, therefore, when performing a \tb-reduction we must
not instantiate the maximal free expressions of the lambda abstraction.
Instead of instantiating them we must substitute a pointer to the single shared
instance in the body of the lambda abstraction. This key idea was first
recognized by Wadsworth [1971]. To illustrate, recall our example from the
previous section

\combinatorbox{}{
    \vspace{-0.5\baselineskip}
    \begin{mlcoded}
        letrec f = g 4 \\
        \phantom{letrec} g = \tlb{x}\tlb{y}+ y (sqrt x) \\
        in \ml{+ (f 1) (f 2)}
    \end{mlcoded}
}

\noindent
The reduction sequence begins in the same way

\vs
\begin{tikzpicture}[
    bullet/.style={circle, fill=black, inner sep=1.5pt},
    arr/.style={-{Stealth[length=2.5mm,width=1.5mm]}, thick},
    node distance=0pt,
    every node/.style={font=\sffamily\footnotesize}
    ]

    \def\rs{0.9}
    \def\as{0.55}

    \pgfmathsetmacro\rowa{0}
    \pgfmathsetmacro\rowb{-1*\rs}
    \pgfmathsetmacro\rowc{-1*\rs-\as}
    \pgfmathsetmacro\rowd{-2*\rs-\as}
    \pgfmathsetmacro\rowe{-2*\rs-2*\as}

    \def\tgtx{4.0}

    % Row 1: + (f 1) (f 2)  [no arrow]
    \node[anchor=north west, inner sep=0pt] at (-1.2,\rowa)
    {\phantom{$\rightarrow$ } \ml{+ (f 1) (f 2)}};

    % Row 2: -> + (• 1) (• 2)   --> ((λx.λy.+ y (sqrt x)) 4)
    \node[anchor=west] at (-1.2,\rowb) {$\rightarrow$ \ml{+ (}};
    \node[bullet] (b2a) at (0.25,\rowb) {};
    \node[anchor=west] at (0.35,\rowb) {\ml{1) (}};
    \node[bullet] (b2b) at (1.15,\rowb) {};
    \node[anchor=west] at (1.25,\rowb) {\ml{2)}};
    \node[anchor=west] (r2tgt) at (\tgtx,\rowc) {\ml{((\tlb{x}\tlb{y}+ y (sqrt x)) 4)}};
    \draw (b2a.south) -- (0.25,\rowc);
    \draw[arr] (0.25,\rowc) -- (r2tgt.west);
    \draw (b2b.south) -- (1.15,\rowc);
    \draw (1.15,\rowc) -- (\tgtx,\rowc);

    % Row 3: -> + (• 1) (• 2)   --> (λy.+ y (sqrt 4))
    \node[anchor=west] at (-1.2,\rowd) {$\rightarrow$ \ml{+ (}};
    \node[bullet] (b3a) at (0.25,\rowd) {};
    \node[anchor=west] at (0.35,\rowd) {\ml{1) (}};
    \node[bullet] (b3b) at (1.15,\rowd) {};
    \node[anchor=west] at (1.25,\rowd) {\ml{2)}};
    \node[anchor=west] (r3tgt) at (\tgtx,\rowe) {\ml{(\tlb{y}+ y (sqrt 4))}};
    \draw (b3a.south) -- (0.25,\rowe);
    \draw[arr] (0.25,\rowe) -- (r3tgt.west);
    \draw (b3b.south) -- (1.15,\rowe);
    \draw (1.15,\rowe) -- (\tgtx,\rowe);

\end{tikzpicture}
\vs

\noindent
But now we see that \ml{(sqrt 4)} is free in the \tl y abstraction, and hence should not
be instantiated when the abstraction is applied. Thus we get

\vs
\begin{tikzpicture}[
    bullet/.style={circle, fill=black, inner sep=1.5pt},
    arr/.style={-{Stealth[length=2.5mm,width=1.5mm]}, thick},
    node distance=0pt,
    every node/.style={font=\sffamily\footnotesize}
    ]

    \def\tgtx{4.0}   % x position of target expression
    \def\rise{0.45}  % y of the top horizontal rail (above code line at y=0)
    \def\drop{0.65}  % y of the bottom horizontal rail (below code line)
    \def\sqrtx{5.75} % x above (sqrt 4) in target

    % Code line at y=0
    \node[anchor=west] at (-1.2,0) {$\rightarrow$ \ml{+ (}};
    \node[bullet] (b1a) at (0.25,0) {};
    \node[anchor=west] at (0.35,0) {\ml{1) (+ 2}};
    \node[bullet] (b1b) at (1.85,0) {};
    \node[anchor=west] at (1.95,0) {\ml{)}};

    % Target expression at y=-\drop
    \node[anchor=west] (tgt) at (\tgtx,-\drop) {\ml{(\tlb{y}+ y (sqrt 4))}};

    % Left bullet: straight south to -\drop, then east with arrowhead to target.west
    \draw (b1a.south) -- (0.25,-\drop);
    \draw[arr] (0.25,-\drop) -- (tgt.west);

    % Right bullet: straight north to +\rise, then east to \sqrtx, then south to
    % just above the text (y=-\drop+0.15 to avoid overlapping)
    \draw (b1b.north) -- (1.85,\rise);
    \draw (1.85,\rise) -- (\sqrtx,\rise);
    \draw[arr] (\sqrtx,\rise) -- (\sqrtx,-\drop+0.18);

\end{tikzpicture}
\vs

\noindent
The instance contains a pointer back to the \ml{(sqrt 4)} in the body of the
abstraction.

\vs
\begin{tikzpicture}[
    bullet/.style={circle, fill=black, inner sep=1.5pt},
    arr/.style={-{Stealth[length=2.5mm,width=1.5mm]}, thick},
    node distance=0pt,
    every node/.style={font=\sffamily\footnotesize}
    ]

    \def\rs{0.9}
    \def\as{0.55}
    \def\tgtx{4.0}
    \def\rise{0.45}
    \def\sqrtx{5.5}

    % Row 1 sits at y=0. Its arrow baseline is just \as below it.
    % Subsequent code rows step by \rs from the arrow baseline of the previous row.
    \pgfmathsetmacro\rowa{0}
    \pgfmathsetmacro\rowb{0-\as}           % arrow baseline for row 1
    \pgfmathsetmacro\rowc{0-\as-\rs}       % row 2 code line
    \pgfmathsetmacro\rowd{0-\as-\rs-\as}   % arrow baseline for row 2
    \pgfmathsetmacro\rowe{0-\as-2*\rs-\as} % row 3 code line
    \pgfmathsetmacro\rowf{0-\as-2*\rs-2*\as} % arrow baseline for row 3 / "2" node

    % Row 1: -> + (• 1) (+ 2 •)
    \node[anchor=west] at (-1.2,\rowa) {$\rightarrow$ \ml{+ (}};
    \node[bullet] (b1a) at (0.25,\rowa) {};
    \node[anchor=west] at (0.35,\rowa) {\ml{1) (+ 2}};
    \node[bullet] (b1b) at (1.85,\rowa) {};
    \node[anchor=west] at (1.95,\rowa) {\ml{)}};
    \node[anchor=west] (r1tgt) at (\tgtx,\rowb) {\ml{(\tlb{y}+ y 2)}};
    \draw (b1a.south) -- (0.25,\rowb);
    \draw[arr] (0.25,\rowb) -- (r1tgt.west);
    \draw (b1b.north) -- (1.85,\rise);
    \draw (1.85,\rise) -- (\sqrtx,\rise);
    \draw[arr] (\sqrtx,\rise) -- (\sqrtx,\rowb+0.18);

    % Row 2: -> + (• 1) 4
    \node[anchor=west] at (-1.2,\rowc) {$\rightarrow$ \ml{+ (}};
    \node[bullet] (b2a) at (0.25,\rowc) {};
    \node[anchor=west] at (0.35,\rowc) {\ml{1) 4}};
    \node[anchor=west] (r2tgt) at (\tgtx,\rowd) {\ml{(\tlb{y}+ y 2)}};
    \draw (b2a.south) -- (0.25,\rowd);
    \draw[arr] (0.25,\rowd) -- (r2tgt.west);

    % Row 3: -> + (+ 1 •) 4
    \node[anchor=west] at (-1.2,\rowe) {$\rightarrow$ \ml{+ (+ 1}};
    \node[bullet] (b3b) at (0.75,\rowe) {};
    \node[anchor=west] at (0.85,\rowe) {\ml{) 4}};
    \pgfmathsetmacro\localrise{\rowe+0.35}
    \draw (b3b.north) -- (0.75,\localrise);
    \draw (0.75,\localrise) -- (\sqrtx,\localrise);
    \draw[arr] (\sqrtx,\localrise) -- (\sqrtx,\rowf+0.18);
    \node[anchor=north] at (\sqrtx,\rowf+0.25) {\ml{2}};

    % Final lines
    \node[anchor=north west, align=left, inner sep=0pt] at (-1.2,\rowf-0.3) {%
        $\rightarrow$ \ml{+ 3 4}\\[2pt]
        $\rightarrow$ \ml{7}%
    };

\end{tikzpicture}
\vs

\noindent
Now the \ml{(sqrt 4)} is only evaluated once, as we had hoped.

\section{Lambda-lifting using Maximal Free Expressions}

In order to achieve full laziness in the lambda reducer of Chapter~12 we
appear to need to identify maximal free expressions dynamically. As we noted
there, this is rather difficult to do efficiently.

Fortunately, it turns out that we can modify the lambda-lifting algorithm so
that a straightforward implementation of the resulting supercombinator
program is automatically fully lazy. The algorithm was invented by Hughes
[1984].

\subsection{Modifying the Lambda-lifting Algorithm}

The modification we need is to \textit{abstract the maximal free expressions}, rather
than free variables, when lambda-lifting a lambda abstraction.

In our running example, the function \ml{g} has the lambda abstraction
\begin{mlcoded}
    \tlb{x}\tlb{y}+ y (sqrt x)
\end{mlcoded}
When doing lambda-lifting on the \ml{\tl y} abstraction, we abstracted \ml{x} out as an
extra parameter, since it occurs free. Instead we should abstract out the entire
(free) subexpression \ml{(sqrt x)} as an extra parameter, thus generating the
supercombinator
\begin{mlcoded}
    \$g1 sqrtx y = + y sqrtx
\end{mlcoded}
The name `\ml{sqrtx}' is an arbitrary name invented for the extra parameter. We
replace the \ml{\tl y} abstraction with the supercombinator \ml{\$g1} applied to the
subexpression, thus:
\begin{mlcoded}
    \tlb{x}\$g1 (sqrt x)
\end{mlcoded}
Completing the compilation in the normal way gives

\combinatorbox{
    \ml{\$g1 sqrtx y = + y sqrtx} \\
    \indent\ml{\$g x = \$g1 (sqrt x)} \\
    \indent\ml{\$f = \$g 4} \\
    \indent\ml{\$Prog = + (\$f 1) (\$f 2)}
}{
    \ml{\$Prog}
}

We get an extra supercombinator because we lose an opportunity for
\ml{\te}-reduction when we take out \ml{(sqrt x)} as an extra parameter rather than \ml{x}.
However, in compensation, the execution will be fully lazy, because the uses
of \ml{(sqrt x)} will be shared. Now we can follow the reduction sequence again:

\vs
\begin{tikzpicture}[
    bullet/.style={circle, fill=black, inner sep=1.5pt},
    arr/.style={-{Stealth[length=2.5mm,width=1.5mm]}, thick},
    node distance=0pt,
    every node/.style={font=\sffamily\footnotesize}
    ]

    \def\rs{0.9}
    \def\as{0.55}
    \def\tgtx{4.2}
    \def\rise{0.45}
    \def\sqrtx{5.25}

    % Row baselines: each code row followed by its arrow baseline
    \pgfmathsetmacro\r{0}          % "$Prog" label
    \pgfmathsetmacro\ra{-1*\rs}    % row 1 code
    \pgfmathsetmacro\raa{-1*\rs-\as} % row 1 arrow baseline
    \pgfmathsetmacro\rb{-2*\rs-\as}  % row 2 code
    \pgfmathsetmacro\rba{-2*\rs-2*\as} % row 2 arrow baseline
    \pgfmathsetmacro\rc{-3*\rs-2*\as}  % row 3 code
    \pgfmathsetmacro\rca{-3*\rs-3*\as} % row 3 arrow baseline
    \pgfmathsetmacro\rd{-4*\rs-3*\as}  % row 4 code
    \pgfmathsetmacro\rda{-4*\rs-4*\as} % row 4 arrow baseline
    \pgfmathsetmacro\re{-5*\rs-4*\as}  % row 5 code
    \pgfmathsetmacro\rea{-5*\rs-5*\as} % row 5 arrow baseline
    \pgfmathsetmacro\rf{-6*\rs-5*\as}  % row 6 code
    \pgfmathsetmacro\rfa{-6*\rs-6*\as} % "2" node level
    \pgfmathsetmacro\rg{-7*\rs-5*\as}  % final lines

    % "$Prog" header
    \node[anchor=north west, inner sep=0pt] at (-1.2,\r)
    {\phantom{$\rightarrow$} \ml{\$Prog}};

    % Row 1: -> + (• 1) (• 2)  -->  ($g 4)
    \node[anchor=west] at (-1.2,\ra) {$\rightarrow$ \ml{+ (}};
    \node[bullet] (b1a) at (0.25,\ra) {};
    \node[anchor=west] at (0.35,\ra) {\ml{1) (}};
    \node[bullet] (b1b) at (1.15,\ra) {};
    \node[anchor=west] at (1.25,\ra) {\ml{2)}};
    \node[anchor=west] (r1tgt) at (\tgtx,\raa) {\ml{(\$g 4)}};
    \draw (b1a.south) -- (0.25,\raa);
    \draw[arr] (0.25,\raa) -- (r1tgt.west);
    \draw (b1b.south) -- (1.15,\raa);
    \draw (1.15,\raa) -- (\tgtx,\raa);

    % Row 2: -> + (• 1) (• 2)  -->  ($g1 (sqrt 4))
    \node[anchor=west] at (-1.2,\rb) {$\rightarrow$ \ml{+ (}};
    \node[bullet] (b2a) at (0.25,\rb) {};
    \node[anchor=west] at (0.35,\rb) {\ml{1) (}};
    \node[bullet] (b2b) at (1.15,\rb) {};
    \node[anchor=west] at (1.25,\rb) {\ml{2)}};
    \node[anchor=west] (r2tgt) at (\tgtx,\rba) {\ml{(\$g1 (sqrt 4))}};
    \draw (b2a.south) -- (0.25,\rba);
    \draw[arr] (0.25,\rba) -- (r2tgt.west);
    \draw (b2b.south) -- (1.15,\rba);
    \draw (1.15,\rba) -- (\tgtx,\rba);

    % Row 3: -> + (• 1) (+ 2 •)  left bullet south-east, right bullet north-east-south onto ($g1 (sqrt 4))
    \node[anchor=west] at (-1.2,\rc) {$\rightarrow$ \ml{+ (}};
    \node[bullet] (b3a) at (0.25,\rc) {};
    \node[anchor=west] at (0.35,\rc) {\ml{1) (+ 2}};
    \node[bullet] (b3b) at (1.85,\rc) {};
    \node[anchor=west] at (1.95,\rc) {\ml{)}};
    \node[anchor=west] (r3tgt) at (\tgtx,\rca) {\ml{(\$g1 (sqrt 4))}};
    \draw (b3a.south) -- (0.25,\rca);
    \draw[arr] (0.25,\rca) -- (r3tgt.west);
    \pgfmathsetmacro\riserowc{\rc+0.35}
    \draw (b3b.north) -- (1.85,\riserowc);
    \draw (1.85,\riserowc) -- (\sqrtx,\riserowc);
    \draw[arr] (\sqrtx,\riserowc) -- (\sqrtx,\rca+0.18);

    % Row 4: -> + (• 1) (+ 2 •)  -->  ($g1 2)
    \node[anchor=west] at (-1.2,\rd) {$\rightarrow$ \ml{+ (}};
    \node[bullet] (b4a) at (0.25,\rd) {};
    \node[anchor=west] at (0.35,\rd) {\ml{1) (+ 2}};
    \node[bullet] (b4b) at (1.85,\rd) {};
    \node[anchor=west] at (1.95,\rd) {\ml{)}};
    \node[anchor=west] (r4tgt) at (\tgtx,\rda) {\ml{(\$g1 2)}};
    \pgfmathsetmacro\riserow{\rd+0.35}
    \draw (b4a.south) -- (0.25,\rda);
    \draw[arr] (0.25,\rda) -- (r4tgt.west);
    \draw (b4b.north) -- (1.85,\riserow);
    \draw (1.85,\riserow) -- (\sqrtx,\riserow);
    \draw[arr] (\sqrtx,\riserow) -- (\sqrtx,\rda+0.18);

    % Row 5: -> + (• 1) 4  -->  ($g1 2)
    \node[anchor=west] at (-1.2,\re) {$\rightarrow$ \ml{+ (}};
    \node[bullet] (b5a) at (0.25,\re) {};
    \node[anchor=west] at (0.35,\re) {\ml{1) 4}};
    \node[anchor=west] (r5tgt) at (\tgtx,\rea) {\ml{(\$g1 2)}};
    \draw (b5a.south) -- (0.25,\rea);
    \draw[arr] (0.25,\rea) -- (r5tgt.west);

    % Row 6: -> + (+ • 1) 4  bullet north-east-south onto "2"
    \node[anchor=west] at (-1.2,\rf) {$\rightarrow$ \ml{+ (+}};
    \node[bullet] (b6b) at (0.6,\rf) {};
    \node[anchor=west] at (0.7,\rf) {\ml{1) 4}};
    \pgfmathsetmacro\riserowf{\rf+0.35}
    \draw (b6b.north) -- (0.6,\riserowf);
    \draw (0.6,\riserowf) -- (\sqrtx,\riserowf);
    \draw[arr] (\sqrtx,\riserowf) -- (\sqrtx,\rfa+0.25);
    \node[anchor=north] at (\sqrtx,\rfa+0.3) {\ml{2}};

    % Final lines
    \node[anchor=north west, align=left, inner sep=0pt] at (-1.2,\rg) {%
        $\rightarrow$ \ml{+ 3 4}\\[2pt]
        $\rightarrow$ \ml{7}%
    };

\end{tikzpicture}
\vs

\noindent
The \ml{(sqrt 4)} is shared, and hence only evaluated once.

So to preserve full laziness we should, during lambda-lifting,
\begin{quote}
    abstract out the maximal free expressions (rather than only the free
    variables) of a lambda abstraction as extra parameters.
\end{quote}
This modification is sufficient to preserve full laziness. We call it \textit{fully lazy
    lambda-lifting}.

One slight optimization is that if a maximal free expression turns out to
have no free variables at all (so it is a CAF), then instead of abstracting it out
as an extra parameter, it can simply be given a name and made into a
supercombinator. The name can then be used instead of the expression. This
is illustrated in the next section.

\subsection{Fully Lazy Lambda-lifting in the Presence of \ml{letrec}s}

As in Chapter~14, our strategy needs to take account of \ml{letrec}s. Consider the
program

\combinatorbox{}{
    \vspace{-0.5\baselineskip}
    \begin{mlcoded}
        \begin{tabular}{l@{}l@{}l}
            let & \\
            &f = \tlb{x} & letrec fac = \tlb{n}(\ldots) \\
            & & in + x (fac 1000) \\
            in &\\
            &+ (f 3) & (f 4)
        \end{tabular}
    \end{mlcoded}
    \vspace{-0.5\baselineskip}
}

The algorithm of Chapter~14 will compile it to

\combinatorbox{
    \begin{mlcoded}
        \$fac fac n = (\ldots) \\
        \$f x = letrec fac = \$fac fac \\
        \phantom{\$f x =} in + x (fac 1000)
    \end{mlcoded}
    \vspace{-\baselineskip}
}{
    \ml{+ (\$f 3) (\$f 4)}
}

\noindent
The function \ml{fac} is defined locally in the body of \ml{f}, and hence \ml{(fac 1000)} cannot
be lifted out as a free expression from the body of \ml{f}. Unfortunately, this means
that \ml{(fac 1000)} will be recomputed each time \ml{\$f} is applied, so we have lost full
laziness.

The solution is to recognize that the definition of \ml{fac} does not depend on \ml{x}.
With this in mind we can `float' the \ml{letrec} for \ml{fac} outwards, giving this program

\combinatorbox{}{
    \ml{letrec}
        \vspace{-0.5\baselineskip}
    \begin{letalign}
        &fac = \tlb{n}(\ldots) \\
        in &let \\
        &f = \tlb{x}+ x (fac 1000) \\
        in \\
        &+ (f 3) (f 4)
    \end{letalign}
    \vspace{-0.5\baselineskip}
}

\noindent
Now our fully lazy lambda-lifter will produce a fully lazy program:

\combinatorbox{
    \begin{mlcoded}
        \$fac n = (\ldots) \\
        \$fac1000 = \$fac 1000 \\
        \$f x = + x \$fac1000 \\
        \$Prog = + (\$f 3) (\$f 4)
    \end{mlcoded}
    \vspace{-0.5\baselineskip}
}{
    \ml{\$Prog}
}

This example also illustrates the utility making a maximal constant
expression into a supercombinator (\ml{\$fac1000} in this case), rather than
abstracting it out as a parameter.

Our strategy now breaks into two phases:
\begin{numbered}
    \item Float out \ml{letrec} (and \ml{let}) definitions as far as possible.
    \item Perform fully lazy lambda-lifting.
\end{numbered}
How far out can a \ml{letrec} definition be floated? The value of a variable bound in
a \ml{letrec} will generally depend on the values of certain free variables. We call
the set of free variables on which a variable \ml{x} depends, \ml{x}'s \textit{free variable set}.
Once we know \ml{x}'s free variable set we can float the definition of \ml{x} outwards
until the next enclosing lambda abstraction binds one of the variables in the
free variable set.

This step has the additional benefit that definitions which have no free
variables at all will be floated out to the top level, where they will be turned
into supercombinators directly.

\section{A Larger Example}

We shall now work through a larger example to show the lambda-lifting
algorithm with full laziness modifications in action.

The example is the function `\ml{foldl}', being used to add up the numbers
between 1 and 100.

\combinatorbox[0.9\textwidth]{
    \begin{mlcoded}
        sumInts n = foldl (+) 0 (count 1 n)
        \\
        \\
        \begin{tabular}{@{}l@{ }l}
            count n m = [\,], &n$>$m \\
            count n m = n : count (n + 1) &m\\
        \end{tabular}
        \\

        \begin{tabular}{@{}l@{ }l}
            foldl op base [\,] &= base \\
            foldl op base (x:xs) &= foldl op (op base x) xs
        \end{tabular}
    \end{mlcoded}
    \vspace{-0.5\baselineskip}
}{
    \ml{sumInts 100}
}

Translating the example into the enriched lambda calculus, we get

\combinatorbox[0.9\textwidth]{}{
    \vspace{-0.5\baselineskip}
    \noindent\ml{letrec}
    \vspace{-0.5\baselineskip}
    \begin{mlcoded}
        sumInts = \tlb{n}foldl (+) 0 (count 1 n) \\
        \\
        count = \tlb{n}\tlb{m}IF ($>$ n m) NIL (CONS n (count (+ n 1) m)) \\
        \\
        foldl\\
%        \vspace{-\baselineskip}
        \begin{tabular}{@{}l@{}l}
                    = \tlb{op}\tlb{base}\tlb{xs} &IF (= xs NIL) \\
            &base \\
            &(foldl op (op base (HEAD xs)) (TAIL xs)) \\
        \end{tabular}
    \end{mlcoded}
    \vspace{-0.5\baselineskip}
    \ml{in sumInts 100}
}

\noindent
(Note: as before, this is not exactly the translation that would be produced by
the pattern-matching compiler, but it suffices for present purposes.) Applying
the algorithm to \ml{foldl}, we choose the innermost lambda abstraction, (\tlb{xs}\ldots),
and look for maximal free expressions, which are \ml{(foldl op)}, \ml{(op base)} and
\ml{base}. We take these out as extra parameters, \ml{p}, \ml{q} and \ml{base} respectively, giving

\combinatorbox[\textwidth]{
    \ml{\$R1 p q base xs = IF (= xs NIL) base (p (q (HEAD xs)) (TAIL xs))}
%    \vspace{-0.5\baselineskip}
}{
    \vspace{-0.5\baselineskip}
    \begin{mlcoded}
        letrec \\
        \phantom{x} sumInts = \tlb{n}foldl + 0 (count 1 n) \\
        \\
        \phantom{x} count = \tlb{n}\tlb{m}IF ($>$ n m) NIL (CONS n (count (+ n 1) m))
        \\
        \\
        \phantom{x} foldl = \tlb{op}\tlb{base}\$R1 (foldl op) (op base) base \\
        in \\
        \phantom{x} sumInts 100
    \end{mlcoded}
    \vspace{-0.5\baselineskip}
}

Now the innermost lambda abstraction is \ml{\tl base}, and its maximal free
expressions are \ml{(\$R1 (foldl op))} and \ml{op}, which we will take out as \ml{r} and \ml{op}
respectively, giving

\combinatorbox[\textwidth]{
    \begin{mlcoded}
        \$R1 p q base xs = IF (= xs NIL) base (p (q (HEAD xs)) (TAIL xs)) \\
        \$R2 r op base = r (op base) base
    \end{mlcoded}
    \vspace{-0.5\baselineskip}
}{
    \vspace{-0.5\baselineskip}
    \begin{mlcoded}
        \ml{letrec} \\
        \phantom{x} sumInts = \tlb{n}foldl + 0 (count 1 n) \\
        \\
        \phantom{x} count = \tlb{n}\tlb{m}IF ($>$ n m) NIL (CONS n (count (+ n 1) m)) \\
        \\
        \phantom{x} foldl = \tlb{op}\$R2 (\$R1 (foldl op)) op \\
        in \\
        \phantom{x} sumInts 100
    \end{mlcoded}
    \vspace{-0.5\baselineskip}
}

Now all the definitions in the top-level \ml{letrec} are supercombinators, because
we have lifted out all the inner lambdas, so after lifting out any constant
expressions we can lift them directly to get

\combinatorbox[\textwidth]{
    \begin{mlcoded}
        \$sumInts n = \$foldlPlus0 (\$count1 n) \\
        \$foldlPlus0 = \$foldl + 0 \\
        \$count1 = \$count 1 \\
        \$count n m = IF ($>$ n m) NIL (CONS n (\$count (+ n 1) m)) \\
        \$foldl op = \$R2 (\$R1 (\$foldl op)) op \\
        \$Prog = \$sumInts 100 \\
        \$R1 p q base xs = IF (= xs NIL) base (p (q (HEAD xs)) (TAIL xs)) \\
        \$R2 r op base = r (op base) base
    \end{mlcoded}
    \vspace{-0.5\baselineskip}
}{
    \ml{\$Prog}
}

\noindent
Notice that we \textit{cannot} eliminate the \ml{op} parameter of \ml{\$foldl}, since it is used
twice on the right-hand side.

\section{Implementing Fully Lazy Lambda-lifting}

We now turn our attention to algorithms for achieving the transformations
required by fully lazy lambda-lifting.

\subsection{Identifying the Maximal Free Expressions}

How can we identify the maximal free expressions of a lambda abstraction?
We can use the concept of \textit{lexical level-number} introduced in Section~13.3.2,
and compute the lexical level of expressions as well as variables. The lexical
level of an expression should be the maximum of the levels of the free
variables within it. Then when lambda-lifting a lambda abstraction at level~$n$,
we should take out as extra parameters any subexpressions within the body
whose level is less than~$n$.

For example, in the base of the lambda abstraction for \ml{g}, which is

\begin{mlcoded}
    \tlb{x}\tlb{y}+ y (sqrt x)
\end{mlcoded}

\noindent
the \ml{\tl x} abstraction is at level 1 and the \ml{\tl y} abstraction is at level 2. Hence the
various subexpressions have level-numbers as follows

\begin{tabular}{l@{\hspace{1.5cm}}l}
    \ml{+}              & level 0 \\
    \ml{(+ y)}          & level 2 \\
    \ml{sqrt}           & level 0 \\
    \ml{(sqrt x)}       & level 1 \\
    \ml{(+ y (sqrt x))} & level 2 \\
\end{tabular}
\vspace{0.5\baselineskip}

\noindent
To summarize:
\vspace{-0.5\baselineskip}
\begin{numbered}
    \item The level-number of a \textit{constant} is 0.
    \item The level-number of a \textit{variable} is the textual nesting depth of the lambda
    which binds it.
    \item The level-number of an \textit{application} \ml{(f x)} is the maximum of the level-numbers of \ml{f} and \ml{x}.
\end{numbered}

\noindent
Given an expression \ml{E}, its \textit{native lambda abstraction} is the enclosing lambda
abstraction whose level-number is the same as that of \ml{E}. Looking `outwards
from \ml{E}' it is the first lambda abstraction which binds any variable in \ml{E}.

All the maximal free expression information can be determined, and
lambda-lifting performed, in a single tree-walk over the expression:
\begin{numbered}
    \item On the way down the tree, the level-number of each lambda abstraction
    is recorded.
    \item On the way up, the level of each expression is computed, using the
    environment and the levels of its subexpressions. If it is applied to
    another expression with the same level-number, then the two are
    merged, otherwise they are given new unique names (since they will be
    maximal free expressions of distinct lambda abstractions). The merging
    is the mechanism whereby free expressions are combined to form
    \textit{maximal} free expressions.

    \item When a lambda is encountered on the way up, it is transformed into a
    supercombinator, and the lambda abstraction is replaced by the super-
    combinator applied to the maximal free expressions. The maximal free
    expressions are those subexpressions with level-number less than that of
    the lambda abstraction, after the merging has taken place.
\end{numbered}

\subsection{Lifting CAFs}

The maximal constant expressions (level 0) need slightly different treatment.
It would be correct to take them out as extra parameters, but there is an easier
way. We can simply define a new supercombinator of zero arguments to be
the constant expression, and use the name of the supercombinator instead of
the expression. No benefit is obtained, however, by doing this with constant
expressions consisting of a single constant (such as \ml{3} or \ml{\$F}), so they can be left
as they are.

For example, in the expression

\begin{mlcoded}
\tlb{x}+ 1 x
\end{mlcoded}

\noindent
the \ml{(+ 1)} is a maximal free expression at level 0, and can be made into a
supercombinator \ml{\$Inc}:

\begin{mlcoded}
\$Inc = + 1
\end{mlcoded}

\noindent
Now the expression becomes

\begin{mlcoded}
\tlb{x}\$Inc x
\end{mlcoded}

\noindent
In this case all that we achieve is the sharing of the \ml{(+ 1)} graph for each
application of the lambda abstraction, but if the constant expression is itself a
redex (like \ml{(+ 1 3)}, for example) then we also save repeated evaluation of the
redex. There was an example of the utility of this in the \ml{\$fac1000} super-
combinator of Section~15.3.2. (Note: there is actually a strong case to be made
for not lifting out a constant expression unless it is in fact a redex -- see Section
15.6.1.)

\subsection{Ordering the Parameters}

In Section~13.3.2 we put the parameters of a supercombinator in order of
increasing level-number, to maximize the opportunities for \ml{\te}-reduction. The
same ordering is useful for maximal free expression parameters, for two
reasons.

The first is the same as before. A maximal free expression will often be just
a single free variable, and in this case, we should still like to have a chance of
\ml{\te}-reduction.

The second reason concerns the \textit{size} of the MFE. To maximize sharing
(which is the object of the exercise) we should like to make our MFEs as large
as possible.

Suppose we have the lambda abstraction

\begin{mlcoded}
    \tlb{x}(\ldots G\ldots F\ldots E\ldots)
\end{mlcoded}

\noindent
where \ml{E}, \ml{F} and \ml{G} are MFEs of the \ml{\tl x} abstraction, and

\begin{mlcoded}
    \normalfont{\normalsize{level of}} F $<$ \normalfont{\normalsize{level of}} G $<$ \normalfont{\normalsize{level of}} E
\end{mlcoded}

\noindent
It would be best to define the supercombinator

\begin{mlcoded}
    \$S f g e x = (\ldots g\ldots f\ldots e\ldots)
\end{mlcoded}

\noindent
and replace the abstraction with

\begin{mlcoded}
    \$S F G E
\end{mlcoded}

\noindent
because then \ml{(\$S F G)} will have a smaller level-number than \ml{E}, and hence will
be taken out of \ml{E}'s native lambda abstraction as a single MFE. If we had
arranged the parameters in the reverse order, \ml{G} and \ml{F} would have had to be
taken out separately.

This will not affect the amount of computation involved (since \ml{(\$S F G)}
cannot be reducible), but it will mean that there is only one instance of the
\ml{(\$S F G)} tree rather than one for each application of \ml{E}'s native lambda
abstraction. Thus, correctly ordering the parameters should make the
maximal free expressions \textit{larger} and \textit{fewer}.

The example in Section~15.4 showed an example of this optimization in
action. We abstracted \ml{(foldl op)}, \ml{(op base)} and \ml{base} from the body of the \ml{\tl x}s
abstraction, calling them \ml{p}, \ml{q} and \ml{base} respectively. Though we did not
mention this at the time, we put \ml{p} first, since \ml{(foldl op)} is freer than \ml{(op base)}
and \ml{base}. This subsequently enabled us to abstract \ml{(\$R1 (foldl op))} from the
\ml{\tl base} abstraction.

We conclude that ordering the parameters by increasing level-number is
better in both these respects.

\subsection{Floating Out the \ml{let}s and \ml{letrec}s}

We recall that to maintain full laziness we must `float' definitions given in \ml{let}s
and \ml{letrec}s outwards. In this section we discuss the algorithm in more detail.

Since we will float out all the definitions in a \ml{letrec} together, we assume that
the dependency analysis described in Chapter~6 has already been performed.
If it were not performed, then a definition might not be floated out as far as
possible, merely because it happened to be defined in the same \ml{letrec} as a
definition which could not be floated out so far. For the same reason we
assume that \ml{let}s contain only a single definition.

In practice, the algorithm of this section could probably be combined with
the dependency analysis algorithm.

How far out should a \ml{let(rec)} be floated? We can compute the `correct'
level-number of its variables, by computing the level-numbers of their
definition bodies. This level-number is correct in the sense that it identifies the
innermost lambda abstraction on which the definition depends. The \ml{let(rec)}
should then be floated out until the nearest enclosing lambda abstraction has
this level-number.

This still leaves some freedom in choosing exactly how far out a \ml{let(rec)} can
be floated. The algorithm which we describe below specifies that:
\begin{numbered}
    \item The immediately enclosing lambda abstraction has the same level-number as that of the variables bound in the \ml{let(rec)}.
    \item The \ml{let(rec)} does not appear in the function position of an application.
    \item It should be floated out as little as possible subject to the constraints (i)
    and (ii).
\end{numbered}

The second condition rejects expressions such as:

\begin{mlcoded}
    (let v = E in E$_1$) E$_2$
\end{mlcoded}

\noindent
in favor of the following equivalent expression, in which the \ml{let} is floated out
one more stage:

\begin{mlcoded}
    let v = E in (E$_1$ E$_2$)
\end{mlcoded}

\noindent
(and similarly for \ml{letrec}s). This has no effect on laziness, but allows an
important simplification in Chapter~20.

The final condition specifies that a \ml{let(rec)} should be floated out no further
than is necessary to meet the first two conditions. To see why this may be
important, consider the expression

\begin{mlcoded}
    IF E (\tlb{x}let v = F in G) H
\end{mlcoded}

\noindent
where \ml{E}, \ml{F}, \ml{G} and \ml{H} are arbitrary expressions, and \ml{F} does not contain \ml{x}. The
algorithm will transform this to

\begin{mlcoded}
    IF E (let v = F in (\tlb{x}G)) H
\end{mlcoded}

\noindent
A sophisticated implementation may be able to avoid constructing the graph
of \ml{H} if \ml{E} turns out to be \ml{TRUE}, and vice versa (see Chapter~20). If we were to
float the \ml{let} out further, we would get the expression

\begin{mlcoded}
    let v = F in (IF E (\tlb{x}G) H)
\end{mlcoded}

\noindent
which is less good, because then the graph of \ml{F} would have to be constructed
whatever value \ml{E} turned out to have.

We can now outline the algorithm as follows. Working from the outside
inwards, for each \ml{let(rec)} perform the following steps:

	\vs
\begin{compactenum}[(1)]
    \item Compute the level-numbers of each definition body. While doing so for
    a \ml{letrec}, assume that the level-number of the variables defined in the
    \ml{letrec} is zero. The reason for this is that the level-number of a recursive
    definition depends only on its free variables, and not on the (as yet
    unknown) level-number of the recursive definition itself.

    \item For a \ml{letrec}, compute the maximum of the level-numbers of the
    definitions' bodies. This is the correct level-number for the variables
    bound in the \ml{letrec}. For \ml{let}s, the correct level-number is that computed
    in Step~1. This level-number should be used for the variables bound in
    the \ml{let(rec)} when processing its body.

    \item Float out the definitions until the next enclosing lambda abstraction has
    the same level-number as that of the variables defined in the \ml{let(rec)},
    which was computed in Step~2.

    \item Finally, if the \ml{let(rec)} now appears in the function position of an
    application, continue to float it out until it does not.
\end{compactenum}
\vs

\noindent
Note: if a \ml{letrec} re-binds a variable that is already in scope, then it cannot be
floated outwards without risk of capturing occurrences of the outer variable.
The solution is to systematically rename one of the variables.

\section{Eliminating Redundant Full Laziness}

The transformations required to achieve full laziness have a price. There are
at least three ways in which we pay:
\begin{numbered}
\item Supercombinators with many arguments (for all the MFEs) are
generated. This increases the size of the redex and slows down reduction.
\item More seriously, more supercombinators may be generated because of
the loss of opportunities for \te{}-optimization. To see this, refer back to the
example in Section~15.4, where three combinators were generated for
\ml{foldl} where one would have sufficed for a non-fully lazy implementation.
More supercombinators mean more reductions.
\item Most serious of all, the program is broken up into small fragments,
fragments of the bodies of functions being exported piecemeal. For a
straightforward template-instantiation implementation this is not a
problem, but if the bodies of supercombinators are compiled then many
opportunities for optimization may be lost. This will become clearer in
Chapter~20, but consider for example the lambda abstraction

\vs
\begin{mlcoded}
\tlb{v}\tlb{x}IF (= v 0) (+ x 1) (+ x 2)
\end{mlcoded}
\vs

The non-fully lazy lambda-lifter will generate a single supercombinator:

\vs
\begin{mlcoded}
    \$R v x = IF (= v 0) (+ x 1) (+ x 2)
\end{mlcoded}
\vs

An optimizing compiler will produce code for \ml{\$R} which first tests the
value of \ml{v}, and then evaluates either the \ml{(+ x 1)} or \ml{(+ x 2)}, to compute a
numerical value. No heap will be consumed. A fully lazy lambda-lifter
will produce two combinators:

\begin{mlcoded}
    \begin{tabular}{@{}l@{$\;=\;$}l}
        \$S1 if-v-zero x & if-v-zero (+ x 1) (+ x 2) \\
        \$S v            & \$S1 (IF (= v 0)) \\
    \end{tabular}
\end{mlcoded}

\noindent
and will replace the \ml{\tl v} abstraction with \ml{\$S}. The compiler will now have to
generate code for \ml{\$S1} to construct \ml{(+ x 1)} \textit{and} \ml{(+ x 2)} \textit{in the heap} before
unwinding the spine of the \ml{if-v-zero} function, about which it now has no
information.
\end{numbered}

\noindent
These objections are substantial, but on the other hand full laziness can save
very large amounts of time and space in some cases. Further study reveals,
however, that the fully lazy lambda-lifter often abstracts out an expression
when \textit{nothing is gained} by so doing. Hence we could improve the trans-
formation by selectively performing ordinary (rather than fully lazy) lambda-
lifting where nothing is gained by the fully lazy method. This section is
therefore devoted to identifying certain situations where fully lazy lambda-
lifting gains nothing, and is based on work by Fairbairn [1985] and Hudak and
Goldberg [1985].

\subsection{Functions Applied to Too Few Arguments}

In the example above, the fully lazy lambda-lifter took out \ml{(IF (= v 0))} as an
extra parameter. However, \ml{IF} requires 3 arguments to reduce, so no work is
saved by sharing this expression. More precisely, just as much work would be
saved by taking out \ml{(= v 0)} as an extra parameter, thus

\begin{mlcoded}
\begin{tabular}{@{}l@{$\;=\;$}l}
    \$T1 v-zero x & IF v-zero (+ x 1) (+ x 2) \\
    \$T v         & \$T1 (= v 0) \\
\end{tabular}
\end{mlcoded}

\noindent
and replacing the \ml{\tl v} abstraction with \ml{\$T}. In a straightforward template-
instantiation implementation some space would be saved by taking out the
larger expression (since the application of \ml{IF} to \ml{(= v 0)} would only be built
once), but even this is not always true in a compiled implementation (see
Chapter~20).

The conclusion is that no work is saved by abstracting out expressions which
consist of a built-in operator or supercombinator applied to too few
arguments. As the example shows, however, the arguments of the function
may be considered for abstraction.

This applies equally to constant expressions which might otherwise be
candidates for a new supercombinator definition (see Section~15.5.2). As an
illustration of this, consider the example in Section~15.4, where the \ml{\$foldlPlus0}
and \ml{\$count1} supercombinators are irreducible; nothing is gained by treating
them as separate supercombinators.

\subsection{Unshared Lambda Abstractions}

Continuing with the same example, suppose the lambda abstraction under
consideration appeared in a context like this:

\begin{mlcoded}
    \begin{tabular}{@{}l}
        let \\
        \quad f = \tlb{v}\tlb{x}IF (= v 0) (+ x 1) (+ x 2) \\
        in \\
        \quad \ldots(f 4 5)\ldots
    \end{tabular}
\end{mlcoded}

\noindent
and suppose that the \ml{(f 4 5)} is the only use of \ml{f}. In this case, the partial
application \ml{(f 4)} \textit{cannot be shared}, since it is used immediately. Using the \ml{\$S}
combinator for \ml{f}, the reduction \ml{(f 4 5)} would go like this:

\begin{letalign}
        f 4 5 & \;\;= \$S 4 5\\
         &$\rightarrow$ \$S1 (= 4 0) 5 \\
         &$\rightarrow$ IF (= 4 0) (+ 5 1) (+ 5 2) \\
         &$\rightarrow$ + 5 2 \\
         &$\rightarrow$ 7
\end{letalign}

\noindent
Since the partial application cannot be shared, neither can the painstakingly
abstracted expression \ml{(= 4 0)}. No sharing would be lost by using the original
\ml{\$R} combinator instead. From this example we can derive a general rule:

\begin{quote}
    given a lambda abstraction \ml{\tlb{x}E} in a context in which it cannot be shared,
    we should not abstract free expressions from \ml{E} because they will not be
    shared. Instead we should abstract only the free variables.
\end{quote}

\noindent
We can justify this rule by observing that free expressions abstracted from \ml{E}
cannot be shared because:
\begin{numbered}
    \item they are not shared inside \ml{E}, since they are abstracted from a single place
    in \ml{E};
    \item they are not shared outside \ml{E}, because the whole lambda abstraction \tlb{x}E
    is not shared.
\end{numbered}

\noindent
The sharing of partial applications is just a specific instance of this general
rule. Notice that for the first time our lambda-lifting strategy becomes \textit{context-dependent}. The trick is to work out when a lambda abstraction might be
shared. This is not at all obvious. For a start, it might be passed as an argument
to another function, in which case a complete analysis would involve looking
at the body of that function. More subtly, consider an extension of our
example:

\begin{mlalign}
        let \\
        & f = \tlb{v}\tlb{x}IF (= v 0) (+ x 1) (+ x 2) \\
        & g = \tlb{x}\tlb{y}+ 1 (f x y) \\
        in \\
        & \ldots \normalsize{\normalfont{expression not mentioning}} f\ldots
\end{mlalign}

\noindent
It does not look as if a partial application of \ml{f} can be shared. But if a partial
application of \ml{g} is shared we will abstract \ml{(f x)} as an MFE from the \ml{\tl y}
abstraction in \ml{g}, so then the partial application of \ml{f} is shared.

Discovering information about sharing is potentially very difficult (it seems
to be another application of \textit{abstract interpretation}; see Chapter~22), but the
saving grace is that we can give up at any time and assume that a partial
application may be shared. The details are beyond the scope of this book but
Fairbairn [1985] and Hudak and Goldberg [1985] each describe their
algorithms.

\section*{References}

\begin{references}

    \item Fairbairn, J. 1985. The design and implementation of a simple typed language based on
    the lambda calculus, pp. 59--60. PhD thesis, \textit{Technical Report 75}. University of
    Cambridge. May.

    \item Hudak, P., and Goldberg, B. 1985. Serial combinators. In \textit{Conference on Functional
        Programming Languages and Computer Architecture, Nancy}, pp. 382--99.
    Jouannaud (editor). LNCS 201. Springer Verlag.

    \item Hughes, R.J.M. 1984. The design and implementation of programming languages.
    PhD thesis, PRG-40. Programming Research Group, Oxford. September.

    \item Wadsworth, C.P. 1971. Semantics and pragmatics of the lambda calculus, Chapter~4.
    PhD thesis, Oxford.

\end{references}